\scanpage{126}
% archive foliation: 125
Soit $P$ une droite projective sur $\mathbb C$. Un ouvert non vide $D\subset P^{\mathrm{an}}$ est dit ``\uline{homogène}'' si le \inserted{sous-}groupe $G$ de $\operatorname{Aut}(P)\xrightarrow{\sim}\operatorname{Aut}(P^{\mathrm{an}})$ $(\simeq\mathrm{GP}(1)_{\mathbb C})$ agit homogènement sur $D$. Si $D$ est homogène et $g\in\operatorname{Aut}(P)$, $g(D)$ est homogène et a comme groupe de stabilité $\operatorname{int}(g).G$.

Le ``théorème d'uniformisation'' nous donne

\uline{Th. (classique) 1.}

a) Il y a exactement trois classes (à op. près de $\operatorname{Aut}(P)$) de domaines homogènes dans $P$, savoir
\[
\begin{aligned}
1^\circ)\quad&D=P\\
2^\circ)\quad&D=P-\{x\}\qquad\text{où }x\in P
\end{aligned}
\]
$3^\circ)$ $D$ un disque ouvert [pour une métrique \inserted{riemannienne} (invariante) sur $X$ --- NB les métriques \inserted{invariantes} (\deleted{sont} toutes multiples scalaires de l'une d'elles, et les disques ouverts pour \deleted{une}\inserted{telle} métrique sont \unsure{indépendants} du choix de celle-ci]

b) Pour $D$ domaine homogène dans $X$, $G\subset\operatorname{Aut}(P)$ son stabilisateur, l'application $G\longrightarrow\operatorname{Aut}(D)$ est un iso. De plus, $D$ est connexe et simplement connexe.

c) Le foncteur $(P,D)\longmapsto D$, de la catégorie \inserted{(pour iso)} des couples $(P,D)$, $P$ droite projective et $D$ domaine homogène de $P$, dans la catégorie \inserted{(pour iso)} des espaces $\mathbb C$-analytiques réguliers de dim 1 connexes et simplement connexes, est une équivalence de catégories.

\uline{Corollaire} Les groupes $G$ de stabilité d'un domaine homogène $D$ de $P$ sont respectivement
\[
\begin{aligned}
1)\quad&\operatorname{Aut}(P)\quad(\simeq\mathrm{GP}(1,\mathbb C))\\
2)\quad&\operatorname{Aut}(\mathbb E^1)\quad(\simeq\operatorname{Aff}(1,\mathbb C))\\
3)\quad&\simeq\mathrm{GP}(1,\mathbb R)\simeq\mathrm{GL}(2,\mathbb R)/\mathbb R^*\simeq\mathrm{SL}(2,\mathbb R)/\mu_2(\mathbb R)
\end{aligned}
\]
(On notera la différence entre $1^\circ)$ et $2^\circ)$ d'une part, où on trouve (des groupes analytiques associés à) un groupe algébrique sur $\mathbb C$, et le cas $3^\circ$
